% GENERATED FILE. Edit canonical lessons, then run npm run generate:books.
% canonical-source-hash: fnv1a64:548d79d6adce3615
% canonical-lessons: HI-C252-bori, HI-C252-dali, HI-C252-kali, HI-C252-kata, HI-C252-kicar

\chapter{A sack, A branch, A bud, A thorn, Mud}
\label{ch:hi-a-sack-a-branch-a-bud-a-thorn-mud}
\hlchaptermodality{252}

\begin{chapteropening}
\textbf{By the end of this chapter:} \emph{I can say a sack, a branch, a bud, a thorn and mud.}

Complete the last lesson of chapter 252: i can say a sack, a branch, a bud, a thorn and mud.
\end{chapteropening}

\section[\texorpdfstring{\hi{बोरी} (borī)}{बोरी (borī)}]{\hi{बोरी} (borī) — a sack}

\label{lesson:HI-C252-bori}



\begin{quote}
Before the new one: say the Hindi for a litre, then the Hindi for a metre.
\end{quote}

\subsection*{You'll want to know: \hi{बोरी}}
\textbf{\hi{बोरी}} — \emph{borī} — \textquotedblleft{}a sack\textquotedblright{}.

\begin{grammarlens}[title={What you've built}]
One more word for this chapter.
\end{grammarlens}

\subsection*{Your turn}
\begin{itemize}
\raggedright
  \item \emph{Say it:} \emph{borī}
  \item \emph{Say it:} \emph{borī}, once more
  \item \emph{Say it:} say \emph{borī}
  \item \emph{Recall:} say the Hindi for a litre, then the Hindi for a metre, then say \emph{borī} again
\end{itemize}

\begin{tcolorbox}[breakable,colback=teal!4,colframe=teal!35!black,title={Before you move on}]
What does \hi{बोरी} mean? (\textquotedblleft{}A sack\textquotedblright{}.) Say it once more.
\end{tcolorbox}
\section[\texorpdfstring{\hi{डाली} (ḍālī)}{डाली (ḍālī)}]{\hi{डाली} (ḍālī) — a branch}

\label{lesson:HI-C252-dali}



\begin{quote}
Before the new one: say the Hindi for a metre, then the Hindi for a sack.
\end{quote}

\subsection*{You'll want to know: \hi{डाली}}
\textbf{\hi{डाली}} — \emph{ḍālī} — \textquotedblleft{}a branch\textquotedblright{}.

\begin{grammarlens}[title={What you've built}]
One more word for this chapter.
\end{grammarlens}

\subsection*{Your turn}
\begin{itemize}
\raggedright
  \item \emph{Say it:} \emph{ḍālī}
  \item \emph{Say it:} \emph{ḍālī}, once more
  \item \emph{Say it:} say \emph{ḍālī}
  \item \emph{Recall:} say the Hindi for a metre, then the Hindi for a sack, then say \emph{ḍālī} again
\end{itemize}

\begin{tcolorbox}[breakable,colback=teal!4,colframe=teal!35!black,title={Before you move on}]
What does \hi{डाली} mean? (\textquotedblleft{}A branch\textquotedblright{}.) Say it once more.
\end{tcolorbox}
\section[\texorpdfstring{\hi{कली} (kalī)}{कली (kalī)}]{\hi{कली} (kalī) — a bud}

\label{lesson:HI-C252-kali}



\begin{quote}
Before the new one: say the Hindi for a sack, then the Hindi for a branch.
\end{quote}

\subsection*{You'll want to know: \hi{कली}}
\textbf{\hi{कली}} — \emph{kalī} — \textquotedblleft{}a bud\textquotedblright{}.

\begin{grammarlens}[title={What you've built}]
One more word for this chapter.
\end{grammarlens}

\subsection*{Your turn}
\begin{itemize}
\raggedright
  \item \emph{Say it:} \emph{kalī}
  \item \emph{Say it:} \emph{kalī}, once more
  \item \emph{Say it:} say \emph{kalī}
  \item \emph{Recall:} say the Hindi for a sack, then the Hindi for a branch, then say \emph{kalī} again
\end{itemize}

\begin{tcolorbox}[breakable,colback=teal!4,colframe=teal!35!black,title={Before you move on}]
What does \hi{कली} mean? (\textquotedblleft{}A bud\textquotedblright{}.) Say it once more.
\end{tcolorbox}
\section[\texorpdfstring{\hi{काँटा} (kā̃ṭā)}{काँटा (kā̃ṭā)}]{\hi{काँटा} (kā̃ṭā) — a thorn, a fork}

\label{lesson:HI-C252-kata}



\begin{quote}
Before the new one: say the Hindi for a branch, then the Hindi for a bud.
\end{quote}

\subsection*{You'll want to know: \hi{काँटा}}
\textbf{\hi{काँटा}} — \emph{kā̃ṭā} — \textquotedblleft{}a thorn, a fork\textquotedblright{}.

\begin{grammarlens}[title={What you've built}]
One more word for this chapter.
\end{grammarlens}

\subsection*{Your turn}
\begin{itemize}
\raggedright
  \item \emph{Say it:} \emph{kā̃ṭā}
  \item \emph{Say it:} \emph{kā̃ṭā}, once more
  \item \emph{Say it:} say \emph{kā̃ṭā}
  \item \emph{Recall:} say the Hindi for a branch, then the Hindi for a bud, then say \emph{kā̃ṭā} again
\end{itemize}

\begin{tcolorbox}[breakable,colback=teal!4,colframe=teal!35!black,title={Before you move on}]
What does \hi{काँटा} mean? (\textquotedblleft{}A thorn, a fork\textquotedblright{}.) Say it once more.
\end{tcolorbox}
\section[\texorpdfstring{\hi{कीचड़} (kīcaṛ)}{कीचड़ (kīcaṛ)}]{\hi{कीचड़} (kīcaṛ) — mud}

\label{lesson:HI-C252-kicar}



\begin{quote}
Before the new one: say the Hindi for a bud, then the Hindi for a thorn.
\end{quote}

\subsection*{You'll want to know: \hi{कीचड़}}
\textbf{\hi{कीचड़}} — \emph{kīcaṛ} — \textquotedblleft{}mud\textquotedblright{}.

\begin{grammarlens}[title={What you've built}]
One more word for this chapter.
\end{grammarlens}

\subsection*{Your turn}
\begin{itemize}
\raggedright
  \item \emph{Say it:} \emph{kīcaṛ}
  \item \emph{Say it:} \emph{kīcaṛ}, once more
  \item \emph{Say it:} say \emph{kīcaṛ}
  \item \emph{Recall:} say the Hindi for a bud, then the Hindi for a thorn, then say \emph{kīcaṛ} again
\end{itemize}

\begin{tcolorbox}[breakable,colback=teal!4,colframe=teal!35!black,title={Before you move on}]
What does \hi{कीचड़} mean? (\textquotedblleft{}Mud\textquotedblright{}.) Say it once more.
\end{tcolorbox}
